\documentclass[../main.tex]{subfiles}
\begin{document}
%==========================================TIÊU ĐỀ TẬP========================================
\begin{table}[h]
    \begin{adjustwidth}{-1cm}{}
        \begin{tabular}{>{\centering\arraybackslash}p{6cm}|>{\raggedright\arraybackslash}p{12.5cm}}
            \multirow{5}{6cm}
            {\\[-40pt]
                \flaregothic\fontsize{20pt}{20pt}\selectfont \centering
                \color{eptype}HISTOIRE\color{black}\\
                \freshbot\fontsize{72pt}{72pt}\selectfont
                \color{epnum}29\\[6pt]
                \cafeteria\fontsize{20pt}{20pt}\selectfont\color{epnum}(D-29)\color{black}
                \\[18pt]
                \color{black}
            }
             &
            {\fotskip\fontsize{18pt}{18pt}\selectfont お\ruby{茶}{ちゃ}\ruby{会}{かい}}
            \\[6pt]
             & {\cafeteria\fontsize{18pt}{18pt}\selectfont Afternoon Tea}                                                                                  \\[4pt]
             & {\vnmsans\fontsize{18pt}{18pt}\selectfont Mùi hương trà chiều}                                                                              \\[8pt]
             & {\caviar{\fontsize{14pt}{14pt}\selectfont Nhân vật: \emp{lamy} \emp{nene} \emp{delu} \emp{kuon2} \emp{suzuno2} \emp{ryuuhou2} \emp{polka}}}
            \\[8pt]
        \end{tabular}
    \end{adjustwidth}
\end{table}
%=========================================TRUYỆN CHÍNH========================================
\color{black}

\txtbx[2]{
    {\color{vol} \uline{Tóm tắt:}} Vài tuần sau những chuyển biến âm thầm tại ký túc xá, các cô gái NePoLaBo bắt đầu ghé thăm Kawachi thường xuyên hơn, và Lamy là người yêu thích không khí yên bình nơi đây nhất. Mang theo những gói trà thảo mộc thơm lành, cô cùng Nene và Suzuno chuẩn bị một buổi trà chiều ấm cúng, có sự góp mặt đầy bất ngờ của Delutaya. Trong khoảnh khắc dạo bước ngang qua hành lang tầng tám, giai điệu dịu êm sau cánh cửa phòng Mikeneko đã níu chân Lamy lại bằng một khúc hát thì thầm và sự thấu cảm sâu sắc về nỗi cô đơn. Một lời tâm sự chân thành với Tổng quản lý Kuon cùng món quà nhỏ để lại bên bậu cửa đã gieo thêm một tia hy vọng dịu ngọt: khi trái tim đã sẵn sàng, một tách trà ấm sẽ luôn đợi chờ để hàn gắn mọi tổn thương.
}

\pov{Yukihana Lamy}

Tháng Sáu gõ cửa Kawachi bằng những vạt nắng đầu hạ óng ả và làn gió mát lành mang theo hơi thở của những vạt rừng xanh ngát phía xa.

Khác hẳn với nhịp sống hối hả, ngột ngạt nơi trung tâm Tokyo, nơi ánh đèn điện phòng thu luôn rực sáng đến tận khuya và tiếng thông báo lịch trình réo rắt liên hồi, vùng ngoại ô này mang một nhịp đập êm đềm đến lạ kỳ. Chỉ cần bước xuống sân ga tàu điện nhỏ bé, hít một hơi thật sâu không khí thoang thoảng mùi cỏ mật và hương đất ẩm, mọi căng thẳng trong lồng ngực tôi dường như tự động tan biến.

Kể từ sau ngày hội ngộ kỳ diệu của \textit{holoFive} vào cuối tháng Ba, ký túc xá nhà Hikawa đã trở thành một chốn về bí mật nhưng vô cùng thân thuộc đối với chúng tôi. Dù tuân thủ nghiêm ngặt những quy ước bảo mật do anh Kuon và chị Kaiho đặt ra, các thành viên NePoLaBo vẫn tìm mọi lý do để ghé qua đây. Botan thì bận rộn với những dự án thể thao điện tử, Polka luôn kín lịch biểu diễn và thu âm, còn Nene-chi thì tràn trề năng lượng nhưng cứ hễ rảnh là lại đòi sang chơi.

Trong số bốn người, tôi có lẽ là người tìm về căn nhà này nhiều nhất.

Ở đây, tôi không cần phải gồng mình khoác lên tấm áo một nàng tiểu thư tuyết thanh lịch trước hàng chục ngàn ánh mắt dõi theo, cũng chẳng phải lo sợ những ánh nhìn soi mói. Tôi thích ngồi bên chiếc bàn gỗ ngoài ban công, lắng nghe tiếng ve sầu râm ran đầu mùa, và tận hưởng cảm giác được là chính mình.

Hôm nay là một buổi chiều thứ Bảy êm ả.

Tôi khẽ xốc lại chiếc túi xách đan bằng cói trên vai. Bên trong là ``kho báu'' mà tôi đã cất công lùng sục suốt cả tuần qua: vài gói trà đen thượng hạng ướp hoa bergamot nhập khẩu, một túi hoa cúc khô do mẹ tôi gửi từ quê lên, và hộp bánh quy nướng bơ hạnh nhân giòn rụm vừa ra lò sáng nay.

``Lamy-chan ơi\~{}! Đợi Nene với nào!''

Một tiếng gọi lanh lảnh quen thuộc vang lên từ sau lưng. Tôi dừng bước, bật cười quay lại.

Nene đang lon ton chạy tới, hai bím tóc màu cam đung đưa theo từng nhịp nhảy chân sáo. Cậu ta diện một chiếc váy yếm màu vàng tươi, hai tay ôm khư khư một chiếc túi giấy to tướng đựng đầy đào tươi vừa mua ở quầy hoa quả trước ga.

``Chạy chậm thôi kẻo ngã bây giờ, Nene-chi,'' tôi dịu dàng nhắc, đưa tay chỉnh lại chiếc nơ cài áo hơi lệch cho cô nàng.

``Tại tui háo hức quá mà! Nghe Lamy-chan bảo hôm nay sẽ tổ chức tiệc trà chiều ở ký túc xá, tui đã phải dậy từ bảy giờ sáng để chọn đào đấy nhé! Đào chín mọng luôn, lát nữa sẽ làm món đào dầm mật ong cho mọi người xem!''

Nhìn nụ cười rạng rỡ như ánh mặt trời ban trưa của Nene, tôi cảm thấy lòng mình nhẹ nhõm đến lạ. Hai chúng tôi sánh bước bên nhau trên con dốc thoai thoải rợp bóng cây, hướng về phía tòa nhà màu xám kiên cố quen thuộc nơi góc phố Kyuukami.

\ct

Cánh cổng lớn mở rộng đón chúng tôi bằng sự bình yên quen thuộc.

Vừa bước qua bậc thềm tầng một, anh Kuon đã xuất hiện ở sảnh chính. Trên tay anh là chiếc kẹp tài liệu quen thuộc, mặc chiếc áo thun đỏ giản dị cùng quần lửng, đôi dép quai hậu phát ra tiếng lẹp xẹp nhẹ nhàng trên nền gạch hoa.

``Chào hai đứa,'' anh Kuon mỉm cười, đôi mắt đỏ thẫm nheo lại vẻ chào đón. ``Đến sớm hơn anh tưởng đấy. Lamy-chan mang theo cả một tiệm trà di động thế kia à?''

``Em chào anh ạ, Kuon-senpai!'' Nene giơ túi đào lên khoe ríu rít. ``Hôm nay NePoLaBo bọn em tới quậy tung ký túc xá đây ạ!''

``Đừng có quậy là anh mừng rồi,'' anh Tổng bật cười, rồi đưa tay đỡ lấy chiếc túi nặng trên tay tôi. ``Suzuno đang ở trên bếp tầng sáu chuẩn bị sẵn ấm chén rồi đấy. Em ấy nghe bảo hai đsa sang chơi nên từ sáng đã lau sạch bóng cả dàn tách sứ cổ của mẹ rồi. Cứ lên tầng năm ở phòng anh mà chờ.''

``Cảm ơn anh. Nọn em xin phép lên trên ạ,'' tôi khẽ gật đầu, lòng dâng lên niềm ấm áp trước sự chu đáo của mọi người.

Căn phòng tầng năm của hai bố con anh Kuon ngập tràn ánh nắng vàng nhạt xiên qua khung cửa sổ kính rộng mở. Trên mặt bàn đảo bằng đá cẩm thạch trắng, Suzuno đã đứng đợi tự bao giờ. Cô bé mặc chiếc váy màu lam trang nhã, bên ngoài khoác chiếc tạp dề trắng viền ren, mái tóc xanh buông nhẹ qua vai. Đôi mắt dị sắc tuyệt đẹp của chị ấy bừng sáng khi thấy chúng tôi bước vào.

``Em chào Lamy-san, chào Nene-san ạ,'' em ấy khẽ cúi đầu chào theo lễ nghi chuẩn mực nhưng giọng nói lại ngập tràn vẻ thân thương. ``Nước tinh khiết đun từ khoáng ngầm đã được căn đúng nhiệt độ rồi ạ. Em cũng đã tráng ấm qua nước sôi như chị dặn lần trước.''

``Suzuno-chan chu đáo quá,'' tôi bước lại gần, cẩn thận mở từng gói trà đặt lên khay gỗ. Mùi hương ngào ngạt của búp trà đen lên men quyện cùng tinh dầu vỏ cam bergamot lập tức lan tỏa, đánh thức mọi giác quan. ``Hôm nay chị mang loại trà bá tước cổ điển này tới, phối thêm một chút hoa cúc khô của em nữa thì vị sẽ ngọt thanh và không bị chát gắt đâu.''

``Tuyệt vời quá ạ!'' Suzuno mỉm cười, đôi gò má ửng hồng vẻ thích thú.

Trong khi tôi và Suzuno đang tỉ mỉ đong từng thìa lá trà vào chiếc ấm gốm men xanh, thì Nene đã không kiên nhẫn nổi mà xắn tay áo lên: ``Để Nene giúp một tay! Pha trà thôi mà, tôi đây cũng là chuyên gia thẩm ẩm thực đấy nhé! Để tui rót nước sôi cho!''

Nói là làm, cậu ấy hăm hở vươn tay chộp lấy chiếc ấm đun nước bằng đồng trên bếp. Thế nhưng chiếc ấm vốn nặng hơn tưởng tượng, lại thêm tay cầm bằng gỗ hơi trơn, vừa nhấc lên khỏi mặt bếp thì cổ tay Nene run bắn lên: ``\textbf{Á á á! Nóng\dots\ trơn quá!}''

Chiếc vòi ấm chao đảo, dòng nước sôi sùng sục suýt chút nữa đã tràn thẳng ra mép bàn và rớt xuống chân của cô nàng.

``\textbf{Cẩn thận!}''

Một bàn tay nhanh như cắt vươn tới, chuẩn xác giữ chặt lấy quai ấm bên trên, trong khi bàn tay còn lại kéo phắt vạt áo Nene lùi về phía sau. Chiếc ấm đun được đặt vững vàng trở lại mặt bếp điện, không một giọt nước nào vương vãi ra sàn.

Tôi giật thót tim, ngẩng đầu lên thì thấy anh Kuon đã đứng cạnh từ lúc nào, hơi thở vẫn giữ được sự bình thản dù vừa ra tay trong chớp mắt.

``Nene-chi à, anh vừa khen xong là hôm nay ngoan ngoãn mà,'' anh thở dài bất lực, đưa tay gõ nhẹ lên trán cô bạn một cái cốc. ``Nước sôi một trăm độ đấy, lóng ngóng một chút là bỏng cả tay bây giờ. Mau ra bàn ngồi ăn bánh đi, để phần pha chế cho Lamy-chan và Suzuno làm.''

Nene ôm trán, xoa xoa chỗ bị gõ, hai má phồng to như hai chiếc bánh bao hấp: ``Hứ! Tại cái ấm này nặng hơn bình thường chứ bộ! Nene chỉ muốn phụ thôi mà\dots''

``Đúng là `nữ hoàng vụng về' của Gen 5 có khác, đi đến đâu là rung chuyển mặt đất đến đấy!'' Một giọng nói giễu cợt đầy quen thuộc vang lên từ phía cửa bếp.

Ryuuhou bước vào, miệng đang ngậm một cây kẹo mút hương dâu, hai tay đút túi chiếc quần jogger rộng thùng thình. Cô bé nhìn chiếc ấm đun nước trên bàn, rồi liếc xéo Nene với ánh mắt đầy trêu chọc: ``May cho bà chị là Kaiho-nee hôm nay đang bận ở xưởng cơ khí cách đây tận bốn phường đấy nhé. Chứ để chị ấy ở nhà mà thấy bà chị suýt làm vỡ cái ấm gốm men ngọc của mẹ thì có mà bị xách cổ ra ban công treo ngược cả buổi chiều!''

``Ryuuhou-chan! Đừng có dọa chị thế chứ!'' Nene kêu ré lên, vội vã chạy vòng ra sau lưng tôi trốn, túm chặt lấy vạt áo tôi như tìm nơi ẩn nấp. Cả gian bếp bỗng chốc vang lên những tràng cười giòn giã. Tôi không kìm được cũng bật cười theo, xoa đầu Nene để dỗ dành cô bạn, trong khi Suzuno nhẹ nhàng đặt chiếc khăn lau tay vào khay.

Bầu không khí rộn ràng ấy cứ thế sưởi ấm căn phòng, xua tan đi bất kỳ khoảng cách nào.

Và rồi, giữa những tiếng cười đùa rộn rã, một âm thanh rất khẽ nơi chiếu nghỉ cầu thang bất chợt lọt vào tai tôi.

*LẠCH CẠCH*

Tiếng dép đi trong nhà nhẹ tênh ngập ngừng dừng lại ở khúc quanh.

Tôi quay đầu nhìn ra hành lang.

Đứng ở mép cửa là một bóng hình nhỏ nhắn trong chiếc áo hoodie màu tím than quá khổ. Mái tóc màu xanh lá hơi rối xõa ngang vai, đôi mắt to tròn long lanh đang nhìn vào gian phòng với vẻ dè dặt, nửa như tò mò, nửa như e ngại.

\uline{Delutaya.}

Tiếng cười của Nene bỗng chốc dịu lại. Ryuuhou cũng dừng màn trêu chọc, khẽ rút cây kẹo mút ra khỏi miệng.

Kể từ sau ngày chúng tôi tụ họp hoom đó, Delu tuy đã cởi mở hơn rất nhiều với gia đình Hikawa, nhưng đây là lần đầu tiên cậu ấy chủ động bước lên tầng năm khi biết có người ngoài ghé thăm. Thói quen thu mình sau hai năm dài ẩn dật vẫn là một lớp vỏ ốc chưa thể biến mất trong một sớm một chiều.

Tôi hít một hơi thật nhẹ, nở nụ cười dịu dàng nhất có thể, cất giọng gọi tự nhiên như thể cậu ấy chưa từng rời xa: ``Delu-chan! Bà xuống đúng lúc lắm, trà vừa ngấm xong này.''

Delutaya khẽ giật mình, ngón tay bấu chặt vào gấu áo hoodie. Cô nàng tóc xanh nhìn tôi, rồi nhìn sang Nene đang vẫy tay rối rít, ánh mắt dần mềm lại.

``Tôi\dots\ ngửi thấy mùi trà thơm quá,'' Delu lí nhí, giọng nói nhỏ nhẹ như tiếng gió thoảng. ``Với lại nghe tiếng Nene-chi la hét từ tận tầng trên\dots''

``A! Delu chê tui ồn ào kìa!'' Nene chu mỏ kêu lên, nhưng ngay sau đó đã lon ton chạy tới kéo tay Delutaya dắt vào bàn. ``Lại đây ngồi với bọn này đi! Lamy-chan pha trà bá tước thơm lừng luôn, tui còn mang cả đào ngọt lịm sang nữa này!''

Bị kéo vào bàn một cách đầy nhiệt tình, Delu không hề phản kháng. Cô bạn ngồi xuống chiếc ghế bọc nệm cạnh tôi, hai bàn tay đan vào nhau đặt trên đùi.

Tôi nâng chiếc ấm sứ men xanh, cẩn thận rót dòng nước trà màu hổ phách óng ả vào chiếc tách sứ trắng tinh trước mặt em. Từng làn khói mỏng bốc lên, mang theo hương bergamot nồng nàn hòa cùng vị ngọt thanh khiết của hoa cúc khô. Tôi đẩy chiếc đĩa nhỏ đựng hai lát bánh quy hạnh nhân sang cạnh tách trà của Delu: ``Loại trà này tôi chọn riêng, có hàm lượng caffeine thấp, phối thêm thảo mộc an thần giúp dễ ngủ đấy. Dạo này bà sáng tác nhạc khuya nhiều, uống cái này sẽ thấy nhẹ đầu hơn.''

Delutaya ngước nhìn tôi, đôi mắt long lanh ngấn nước trong giây lát. Cô nàng nâng tách trà bằng cả hai tay, áp lòng bàn tay vào thành gốm ấm áp, rồi khẽ nhấp một ngụm nhỏ.

Vị ngọt bùi lan tỏa nơi đầu lưỡi. Đôi vai hơi căng cứng của cô bé từ từ buông lỏng, và trên đôi môi nhỏ nhắn, một nụ cười ngượng ngùng nhưng rạng rỡ như đóa hoa đầu mùa khẽ hé nở: ``Ngon lắm\dots\ Ấm áp lắm thật đấy\~{}\dots Cảm ơn Lamy-chan nhiều.''

Nhìn nụ cười ấy, tim tôi như thắt lại vì xúc động.

Không cần những lời hoa mỹ, không cần những hồi ức đau buồn được xới lên. Chỉ cần một tách trà nóng, một chiếc bánh quy và sự hiện diện chân thành bên nhau, chúng tôi đã lại là \textit{holoFive}, trọn vẹn và bình yên như chưa từng có cuộc chia ly nào.

\ct

Buổi trà chiều trôi qua trong những câu chuyện ríu rít không hồi kết.

Nene say sưa kể về những trò nghịch ngợm trong buổi livestream tuần trước, Delutaya thỉnh thoảng lại bật cười khúc khích phụ họa, còn Ryuuhou thì tranh thủ bóc mẽ mấy tật xấu của anh trai khiến anh Kuon chỉ biết lắc đầu cười trừ.

Đến khoảng bốn giờ chiều, Nene kéo Ryuuhou và Delu sang phòng đa năng tầng hai để cùng xem lại mấy video vũ đạo mới. Suzuno ở lại thu dọn ấm chén, còn tôi xin phép mang một bình giữ nhiệt nhỏ đựng nước sôi và bộ dụng cụ pha trà lên dạo quanh tòa nhà một chút để hít thở không khí.

Tôi bước lên cầu thang bộ. Từng bậc thang gỗ lim kêu cọt kẹt êm ái dưới chân.

Khi bước chân tôi chạm đến tầng tám, không gian xung quanh bỗng chốc lắng đọng lại.

Tầng tám là tầng cao nhất của tòa ký túc xá, nằm ngay sát dưới sân thượng. Khác với sự ồn ào, náo nhiệt của tầng năm hay tiếng đàn hát ở tầng hai, nơi đây vắng lặng đến mức tôi có thể nghe rõ tiếng gió chiều rít khẽ qua khe cửa sổ kính ở cuối hành lang.

Ánh hoàng hôn đỏ ối bắt đầu tràn vào, nhuốm những vệt sáng màu cam sẫm lên sàn gỗ dài hun hút.

Và ở tận cùng của hành lang ấy là căn phòng mang số 801.

Căn phòng của \uline{chị Mikeneko}.

Tôi bước chậm lại, từng bước chân nhẹ như mèo bước để không tạo ra tiếng động. Khi chỉ còn cách cánh cửa gỗ chừng vài mét, bước chân tôi đột ngột khựng lại.

Từ sau cánh cửa đóng kín mít, một âm thanh dìu dịu đang phát ra.

Đó không phải là sự im lìm ngột ngạt hay tiếng khóc nức nở mà tôi từng lo sợ. Đó là một bản nhạc lofi không lời, êm dịu và trầm buồn, với tiếng piano lặp đi lặp lại nhịp nhàng trên nền tiếng mưa rơi nhân tạo lách tách.

Tôi đứng lặng người trước cánh cửa.

Bất giác, ký ức về những ngày tháng cũ ùa về trong tâm trí tôi.

Cái tên Uruha Rushia đối với tôi, và với tất cả các thế hệ đàn em, từng là một tiền bối dịu dàng, người luôn sẵn sàng lắng nghe và cổ vũ chúng tôi trong những ngày đầu chập chững bước vào công ty. Nhưng tôi cũng biết, đằng sau sự ngọt ngào ấy là một tâm hồn mỏng manh, dễ tổn thương vô cùng, luôn mang trong mình nỗi sợ hãi tột cùng bị mọi người bỏ rơi.

Và rồi cơn bão nghiệt ngã hồi tháng Hai ập đến. Nó quét sạch tất cả, xé toạc mọi thứ, đẩy chị  ấy xuống vực sâu của sự phán xét và cô độc, trước khi được anh Kuon đưa về che chở trong căn phòng này.

Hôm trước, anh Kuon đã dặn dò chúng tôi rất kỹ: chị Mikeneko vẫn chưa sẵn sàng. Vết thương lòng còn quá mới, niềm tin vào thế giới bên ngoài vẫn còn rạn nứt sâu sắc. Bất kỳ một sự tiếp cận đường đột nào cũng có thể khiến chị ấy hoảng loạn mà thu mình sâu hơn vào bóng tối.

Tôi nhìn cánh cửa gỗ im lìm. Nắm đấm cửa bằng đồng lạnh ngắt.

Tôi không gõ cửa.

Tôi khẽ thở ra một hơi dài, rồi chậm rãi ngồi xuống, tựa lưng vào bức tường đối diện cánh cửa phòng.

Bản nhạc lofi bên trong chuyển sang một giai điệu rất quen, đó là một khúc nhạc nhẹ mà chúng tôi từng nghe chung trong phòng chờ văn phòng hồi năm 2020, cái thuở mà mọi thứ vẫn còn vẹn nguyên và trong sáng.

Giai điệu du dương ấy chạm vào một góc khuất sâu thẳm trong tim tôi.

Tôi khép hờ mắt, ngước nhìn lên trần nhà ngập tràn ánh hoàng hôn, rồi vô thức cất tiếng ngân nga theo điệu nhạc.

Tôi không hát thành lời. Tôi chỉ dùng chất giọng trong trẻo, êm dịu như làn tuyết đầu mùa của mình để hòa vào từng nốt piano đang vang vọng sau cánh cửa. Không có kỹ thuật thanh nhạc phức tạp, không có áp lực biểu diễn, chỉ thuần túy là một khúc hát ru của một người bạn muốn gửi gắm sự vỗ về qua bức tường ngăn cách.

Tôi ngân nga hết một đoạn điệp khúc.

Và rồi, một điều kỳ diệu đã xảy ra.

Bên trong căn phòng, tiếng nhạc piano bỗng chốc nhỏ dần đi.

Nó không hề bị tắt phụt đi như khi người ta khó chịu muốn xua đuổi kẻ làm phiền. Tiếng nhạc chỉ giảm xuống chừng hai nấc âm lượng, vừa đủ nhỏ để tiếng đàn không lấn át, và vừa đủ để người bên trong có thể lắng nghe rõ hơn từng thanh âm thì thầm ngoài hành lang.

Tôi sững người, tim đập thình thịch trong lồng ngực.

Cánh cửa vẫn không hề mở ra. Tay nắm cửa vẫn bất động.

Nhưng sự điều chỉnh âm lượng nhỏ nhoi ấy\dots\ đối với tôi, còn quý giá hơn ngàn vạn lời đáp lại.

Chị ấy đang lắng nghe.

Chị ấy biết có một người đang ngồi ngoài này, không xoi mói, không phán xét, chỉ lặng lẽ hát cùng chị qua một khúc nhạc chiều.

Tôi mỉm cười, khóe mắt chợt cay cay. Tôi tiếp tục ngân nga nốt đoạn nhạc còn lại với tất cả sự dịu dàng của trái tim mình, trước khi bản nhạc bên trong dần đi đến hồi kết trong tiếng thở dài êm ả của buổi hoàng hôn.

\ct

Khi tôi bước xuống tầng năm, ánh chiều tà đã ngả sang màu tím thẫm nơi chân trời.

Anh Kuon đang đứng một mình ngoài ban công lớn, hai tay tựa vào lan can sắt, trên tay là cuốn sổ tay công việc đã gập lại. Gió chiều thổi tung làn tóc đen của anh, để lộ vầng trán mang nét trầm tư quen thuộc của người gánh vác cả gia đình này.

Tôi bước lại gần, đặt chiếc tách sứ bốc khói nghi ngút lên chiếc bàn trà nhỏ bên cạnh: ``Trà nóng đây ạ, Kuon-senpai.''

Anh Tổng giật mình quay lại, nhìn thấy tôi thì nét mặt nghiêm nghị bỗng chốc giãn ra, nở nụ cười hiền: ``Cảm ơn Lamy-chan nhé. Nene-chan với Ryuuhou vẫn đang quậy dưới phòng tập à?''

``Vâng, ba đứa đang thi nhảy bài mới vui lắm ạ,'' tôi kéo chiếc ghế mây ngồi xuống bên cạnh anh, hai bàn tay áp vào tách trà của mình để tìm kiếm hơi ấm. ``Kuon-senpai\dots\ lúc nãy em vừa lên tầng tám.''

Bàn tay đang cầm tách trà của anh Kuon hơi khựng lại. Anh đưa mắt nhìn tôi, ánh mắt thoáng nét lo lắng: ``Em lên tầng tám à? Mike-chan có sao không?''

``Em không gõ cửa đâu ạ,'' tôi vội lắc đầu xua tay để anh yên tâm. ``Em chỉ đứng ngoài hành lang thôi. Em nghe thấy tiếng nhạc trong phòng em ấy phát ra\dots\ rồi em ngồi xuống trước cửa, khẽ hát theo một đoạn.''

Tôi kể lại cho anh nghe khoảnh khắc âm lượng bên trong nhỏ dần đi.

Nghe xong, anh Kuon thở phào một hơi dài, bờ vai dường như nhẹ nhõm đi một chút. Anh nhấp một ngụm trà, ánh mắt nhìn xa xăm về phía những dãy núi trập trùng: ``Thế là tốt rồi\dots\ Đêm hôm trước, em ấy cũng vừa lén đẩy một đĩa bánh quy tự làm qua khe cửa cho anh. Em ấy đang cựa mình thức giấc, từng chút, từng chút một.''

Tôi im lặng nhìn ráng chiều đang tắt dần, lòng dâng lên muôn vàn cảm xúc ngổn ngang: ``Anh biết không, senpai\dots\ Em hiểu cảm giác của em ấy.''

Anh Kuon quay sang nhìn tôi, chăm chú lắng nghe.

``Em hiểu nỗi cô đơn tột cùng đó,'' giọng tôi chùng xuống, mang theo những ký ức sâu kín mà tôi hiếm khi thổ lộ với ai. ``Hồi mùa thu năm 2020\dots\ cái ngày mà Aloe-chan rời đi\dots\ em đã từng sợ hãi đến phát khóc. Em sợ phải bước lên sân khấu một mình. Em sợ ánh mắt của công chúng. Có những đêm, em tự nhốt mình trong phòng kín, trùm chăn khóc nức nở vì sợ rằng chỉ cần chớp mắt một cái, những người đồng đội còn lại cũng sẽ biến mất khỏi cuộc đời mình.''

Tôi siết chặt chiếc tách trà trong tay, cảm nhận hơi nóng ấm áp truyền qua da thịt: ``Nếu hồi đó không có sự kiên trì của Nene, sự vững vàng của Botan, sự lạc quan đến ngốc nghếch của Polka\dots\ và nếu không có anh luôn đứng sau lưng làm chiếc khiên che chắn cho chúng em, có lẽ em đã gục ngã từ lâu rồi. Vì em từng trải qua bóng tối ấy, nên khi nhìn cánh cửa tầng tám đóng chặt, tim em đau lắm. Em chỉ muốn lao vào, ôm lấy em ấy và nói rằng mọi người không hề ghét bỏ hay lãng quên em ấy\dots''

Tôi dừng lại, hít một hơi sâu rồi khẽ cười trừ: ``Nhưng em biết\dots\ em không thể làm thế được. Vết thương của Mikeneko-senpai quá lớn, sự tổn thương ấy không thể lành lặn chỉ bằng vài lời động viên vội vã. Càng ép buộc, em ấy sẽ càng sợ hãi.''

Anh Kuon nhìn tôi thật lâu. Trong ánh mắt đỏ thẫm của người Tổng quản lý, tôi nhìn thấy sự xúc động sâu sắc lẫn niềm tự hào khôn nguôi. Anh đặt nhẹ bàn tay ấm áp lên đầu tôi, khẽ vỗ về như người anh trai lớn trong nhà: ``Em đã trưởng thành nhiều lắm rồi, Lamy-chan.''

``Dạ?'' tôi ngơ ngác chớp mắt.

``Cô ấy hay e ấp, nhút nhát và dễ rơi nước mắt ngày nào của \textit{holoFive}, giờ đã biết thấu hiểu và bao dung cho nỗi đau của người khác đến mức này rồi,'' anh Kuon mỉm cười, giọng nói trầm ấm lan tỏa trong gió chiều. ``Em nói đúng. Sự chữa lành không thể vội vàng. Nó cần thời gian, sự kiên nhẫn, và quan trọng nhất là sự chân thành không đòi hỏi hồi đáp. Nhờ có những người như em, anh tin là một ngày nào đó không xa, cánh cửa tầng tám kia sẽ tự mình mở lối.''

Lời khen giản dị của anh khiến sống mũi tôi cay cay, nhưng trong lòng lại ngập tràn một nguồn năng lượng ấm áp lạ thường.

\ct

Trời sẩm tối, thành phố Kawachi bắt đầu lên đèn.

Đã đến lúc NePoLaBo phải chuẩn bị ra ga để bắt chuyến tàu điện về lại Tokyo. Nene và Delutaya ôm tạm biệt nhau ở tiền sảnh hồi lâu, hẹn nhau cuối tuần sau sẽ lại sang chơi.

Trước khi bước ra khỏi cổng, tôi xin phép anh Kuon năm phút.

Tôi bước nhanh lên tầng tám một lần nữa.

Lần này, trên tay tôi là một gói quà nhỏ nhắn, được bọc cẩn thận bằng giấy gói hoa văn bông tuyết trắng muốt, thắt một chiếc nơ ruy băng màu xanh lam nhạt. Bên trong là gói trà đen ướp hoa bergamot ngon nhất mà tôi mang theo, kèm theo vài thanh đường phèn nhỏ.

Và đính trên chiếc nơ là một tấm thiệp nhỏ hình giọt sương, trên đó tôi nắn nót viết vài dòng chữ mực tím:

\txtbx{
    ``Gửi Mikeneko-san, \\[4pt]
    Trà chiều hôm nay thơm lắm. Gửi tặng chị một chút hương vị mùa hạ này nhé. \\
    Khi nào chị sẵn sàng\dots\ chúng ta cùng uống trà nhé.''\\[4pt]
    \begin{flushright}--- Yukihana Lamy\end{flushright}
}

Tôi cúi người, nhẹ nhàng đặt gói trà lên chiếc khay gỗ nhỏ bên cạnh bậu cửa phòng 801.

Hành lang vẫn im lìm, nhưng ánh đèn vàng ấm áp hắt qua khe cửa sổ như đang lặng lẽ dõi theo từng cử chỉ của tôi. Tôi đứng thẳng người, khẽ cúi đầu chào cánh cửa gỗ, rồi mỉm cười xoay người bước xuống cầu thang.

Tàu điện lăn bánh rời ga Kawachi, mang theo tôi và Nene trở về với Tokyo hoa lệ.

\begin{center}
    \uline{Sáng hôm sau.}
\end{center}

Ánh bình minh đầu ngày rọi những tia nắng vàng rực rỡ qua khung cửa sổ tầng tám.

Anh Kuon trong bộ đồ thể thao đi kiểm tra các tầng theo thói quen buổi sáng. Khi bước chân anh dừng lại trước căn phòng, ánh mắt anh bất chợt chạm vào chiếc khay gỗ nhỏ trước cửa.

Chiếc khay gỗ vẫn nằm nguyên ở đó.

Thế nhưng, gói quà bọc giấy hoa tuyết và tấm thiệp nhỏ nơ xanh\dots\ đã hoàn toàn biến mất.

Thay vào đó, nơi khe cửa hẹp khép hờ, một làn hương thơm thoang thoảng dịu nhẹ bay ra ngoài hành lang: mùi hương của lá trà đen bergamot vừa được pha bằng nước sôi đầu ngày, quyện cùng một tia nắng ấm áp đầu hạ.

Anh Kuon đứng lặng người trong giây lát, rồi khóe môi anh khẽ cong lên một nụ cười rạng rỡ nhất trần đời.

Hạt mầm hy vọng đã được gieo xuống. Và mùa hoa nhất định sẽ nở rộ.

%==========================================TRUYỆN PHỤ=========================================
\omk

\pov{-}

Khoảng năm giờ chiều hôm trước - ngay trong lúc Lamy và Kuon đang đứng tâm sự ngoài ban công tầng năm.

*RẦM!*

Cánh cửa sảnh tầng một bị đẩy bật ra một cách thô bạo, kèm theo một tiếng thở dốc hổn hển như vừa chạy việt dã mười cây số:``Hế\dots\ hế nhô cả nhà! Polka rạp xiếc đã đổ bộ thành công đâaaaay!''

Polka xuất hiện trong bộ trang phục thường ngày sặc sỡ, mái tóc vàng xù xì bù xù vì gió, mồ hôi nhễ nhại chảy ròng ròng trên trán. Cô nàng lảo đảo bước vào phòng khách tầng năm, rồi với một động tác kịch tính như diễn viên xiếc tung hứng, Polka ném một chiếc bao tải nilông to đùng lên mặt bàn trà gỗ: ``Kẹt xe buýt kinh hoàng từ ga Kawachi tới đây! Nhưng không sao! Ta đã mang về chiến lợi phẩm vĩ đại nhất trần đời!''

Chiếc bao tải bung ra, làm tràn ra mặt bàn hàng trăm cây kẹo mút khổng lồ đủ mọi màu sắc sặc sỡ, từ vị dâu tây, bắp rang bơ, coca cho đến vị sầu riêng cay nồng.

Ryuuhou - vừa từ phòng tập nhảy bước ra với chiếc khăn vắt ngang cổ - hai mắt lập tức sáng rực lên như đèn pha ô tô. Cô bé lao vút tới bàn như một mũi tên: ``Oa! Kẹo mút bắp rang bơ phiên bản giới hạn của rạp xiếc! Chia cho em một nửa đi!''

``Không đời nào!'' Polka lập tức dang hai tay ôm khư khư lấy đống kẹo, phồng má trợn mắt: ``Đây là kẹo của Polka mua bằng mồ hôi nước mắt! Muốn ăn thì phải thi đấu!''

``Thi cái gì? Em chấp chị luôn!'' Ryuuhou xắn tay áo lên, tinh thần hiếu chiến trỗi dậy ngút trời.

``Thi ăn kẹo mút siêu tốc! Ai bóc vỏ và ngậm tan hết một cây kẹo mút khổng lồ này trước thì người đó được ôm cả bao tải về!''

``Chơi luôn! Nhào vô!''

Thế là ngay giữa sàn phòng khách, hai con người - một cô nàng Fennec rạp xiếc và một cô bé rocker 16 tuổi - ngồi bệt xuống đất, phồng mang trợn má nhai kẹo lạo xạo như hai chú sóc chuột tranh nhau hạt dẻ. Nene-chi đứng bên cạnh thì nhảy cẫng lên, hai tay vỗ bôm bốp cổ vũ nhiệt tình: ``Cố lên Polka-senpai! Đừng để thua Ryuuhou-chan! Cố lên Ryuuhou-chan, cắn nát cây kẹo ra đi!''

Tiếng hò reo, tiếng nhai kẹo rồm rộp và tiếng la ó ầm ĩ vang dội khắp cả tầng năm, làm rung chuyển cả mấy khung tranh treo tường.

Đúng lúc cuộc thi đang bước vào giai đoạn gay cấn nhất — Polka nghẹn kẹo trợn tròn mắt còn Ryuuhou thì méo cả quai hàm — thì một luồng sát khí lạnh ngắt bất chợt ập tới từ lối cầu thang dẫn lên phòng khách.

*CẠCH.*

Cánh cửa phòng khách mở toang.

Kaiho đứng sững ở ngưỡng cửa.

Cô vừa chạy xe máy từ xưởng cơ khí cách bốn phường về nhà sau ca làm muộn, trên người vẫn mặc chiếc áo khoác lao động màu xám lấm lem vài vệt dầu mỡ, mái tóc đen cột đuôi ngựa gọn gàng. Vừa bước chân lên tầng năm định thay đồ thì đập vào mắt cô là khung cảnh hỗn loạn này. Đôi mắt xanh biếc sắc lẹm như lưỡi dao cạo quét một vòng từ đống vỏ kẹo vương vãi trên sàn đến hai kẻ đang phồng mồm trợn mắt dưới đất.

Cả căn phòng bỗng chốc rơi vào sự im lặng chết chóc.

Nene lập tức ngậm chặt miệng, từ từ lùi lại ba bước rồi núp sau lưng chiếc ghế sofa.

``Hai người\dots'' giọng Kaiho trầm xuống, lạnh buốt như băng tuyết Bắc Cực, ``đang làm cái trò hề gì ở đây thế?''

Polka giật mình nghẹn họng, ho sù sụ: ``Khụ\dots\ khụ! Bọn em\dots\ thi đấu công bằng\dots''

Ryuuhou mặt cắt không còn giọt máu, cây kẹo mút dính chặt vào má: ``Kaiho-nee\dots?''

Kaiho không nói nửa lời. Cô bước thẳng tới bàn, vươn cánh tay rắn rỏi chụp lấy miệng chiếc bao tải kẹo khổng lồ, giật phắt lên không trung nhẹ như không. Bàn tay còn lại vung ra chộp lấy luôn hai que kẹo đang ngậm dở trên miệng của Polka và Ryuuhou rồi vứt thẳng vào sọt rác.

``Tịch thu toàn bộ.''

``\textbf{HẢAAA!?}'' cả Polka lẫn Ryuuhou đồng thanh kêu ré lên.

``Sau năm giờ chiều không được nạp lượng đường vượt quá ba mươi gam. Quy định nội quy ký túc xá, điều số mười bốn,'' Kaiho lạnh lùng tuyên bố, vác chiếc bao tải kẹo nặng trịch lên vai như vác bao cát. ``Còn gây ồn ào nữa thì cắt cơm tối cả hai đứa. Dọn sạch sàn nhà trong mười phút, nghe rõ chưa?''

Nói xong, cô quay lưng sải bước dài về phía cầu thang, bóng lưng thẳng tắp đầy uy quyền không ai dám cãi nửa lời.

Polka và Ryuuhou ngồi bệt dưới sàn, nhìn nhau với đôi mắt đẫm lệ, đồng thanh gào lên trong tuyệt vọng: ``\textbf{Kaiho-nee/Kaigou-san là đồ ác quỷ!!!}''

Ở góc phòng khách, Naomi vừa ôm chú gấu bông bước ra, nhìn cảnh tượng ấy rồi ngơ ngác hỏi Nene: ``Ủa chị Nene\dots\ hai chị kia đang tập diễn kịch bi kịch ạ?''

\end{document}